Ukraine elected half of the Verkhovna Rada by party list in all six elections
(the whole of it in 2006 and 2007); BLUE\_DB reports the party-list vote.

\paragraph{2012, 2014, 2019: hromadas.}  The Central Election Commission (CVK)
publishes the party-list vote of each polling station, with the station's
address.  Each station is placed in the settlement its address names, matched
against the settlement layer of the OCHA administrative boundaries (COD-AB
v05, 2025), and so in the \textit{hromada} that settlement belongs to today:
the hromadas are Ukraine's municipalities since the 2020 reform, and the three
elections are expressed in this one current geography (\texttt{UA\_0502003},
from the KATOTTH-based code \texttt{UA0502003}).  Settlements named more than
once in an oblast are resolved by the other stations of the same old raion;
settlements renamed since (the 2016 decommunisation, the 2021--24 renamings)
by the hromadas' former names, the same station number in a later election, a
short table of renamed cities, or the postcode.  Crimea, which voted only in
2012, keeps its village councils in the boundary file.  Votes abroad are
\texttt{UA\_UAZZ99}, votes on ships and at the polar station \texttt{UA\_UAYY99};
in 2012 and 2014 the CVK publishes the abroad vote only as a total, so its
party votes are the national result less the stations at home, and its
registered voters and turnout are not available.  In 2012 the stations report
registered voters and valid votes only, so turnout is the valid vote.

\paragraph{2006, 2007: territorial constituencies.}  These elections are
published by territorial constituency (ТВО, \texttt{UA\_T001}--\texttt{UA\_T225}),
each a union of pre-2020 raions, cities of oblast significance and city
districts.  None of these units survives in today's boundaries, so they are
rebuilt: every settlement takes the old raion or city of its polling stations'
addresses in 2012, settlements without a station that of their nearest
neighbour, and each hromada is shared out among its settlements (whole, or in
Voronoi cells around them).  The districts of the eleven cities split between
constituencies come from OpenStreetMap, matched to their 2006 names through
the \texttt{old\_name} tags and Wikidata.  The constituency outlines are
therefore an approximation where old raion borders cut through today's
hromadas.

\paragraph{2002.}  The 2002 party-list vote is published by single-member
constituency only, and 134 of the 225 constituencies are drawn through
raions and city districts by street; it is not yet part of the distribution.

\paragraph{Regions.}  Ukraine is outside NUTS.  BLUE\_DB uses a pseudo-NUTS
hierarchy: NUTS~1 = NUTS~2 = the 27 oblast-level regions (\texttt{UA1}/\texttt{UA11}
\ldots), NUTS~3 = the 2020 raions (\texttt{UA111} \ldots; Kyiv and Sevastopol
are their own raion).  A territorial constituency is placed in the raion
holding most of its area.
